\section{The Interpretability Landscape}

\begin{frame}{Why Explain a Model At All?}

    \begin{columns}[T]
        \begin{column}{0.49\textwidth}
            \begin{itemize}
                \setlength\itemsep{0.55em}
                \item \textbf{Debugging.} By far the most common real use. Explanations expose
                      leakage, shortcut features, and broken preprocessing that accuracy alone
                      hides.
                \item \textbf{Trust and adoption.} A clinician or underwriter who cannot see the
                      reasoning will not use the tool --- or will over-trust it.
                \item \textbf{Recourse.} Telling a rejected applicant \emph{what to change} is a
                      different requirement from telling an engineer what the model used.
            \end{itemize}
        \end{column}
        \begin{column}{0.49\textwidth}
            \begin{itemize}
                \setlength\itemsep{0.55em}
                \item \textbf{Accountability.} When a decision is challenged, somebody has to be
                      able to say why the model produced it. ``The network decided'' is not an
                      answer anyone accepts twice.
                \item \textbf{Science.} In a scientific application the model \emph{is} the
                      hypothesis, and what it uses is the result you wanted.
            \end{itemize}
            \vspace{0.5em}
            \begin{tcolorbox}[colback=gray!8, boxrule=0pt]
                \small Each purpose demands a different artefact. Decide the purpose \textbf{before}
                choosing the method.
            \end{tcolorbox}
        \end{column}
    \end{columns}
\end{frame}

\begin{frame}{Three Axes for Sorting Every Method}

    \centering
    \small
    \begin{tabular}{p{2.9cm} p{4.4cm} p{4.4cm}}
        \toprule
        \textbf{Axis} & & \\
        \midrule
        \textbf{Intrinsic} vs \textbf{post-hoc}
        & The model is readable as it stands: linear/logistic regression, a shallow tree, a
          rule list, a GAM.
        & An external procedure explains a fitted black box: permutation importance, SHAP, LIME,
          Grad-CAM. \\
        \addlinespace
        \textbf{Global} vs \textbf{local}
        & What does the model do \emph{in general}? Feature importances, partial dependence.
        & Why \emph{this} prediction for \emph{this} row? SHAP values, LIME, counterfactuals. \\
        \addlinespace
        \textbf{Model-specific} vs \textbf{model-agnostic}
        & Uses the internals: tree paths, gradients, attention, TreeSHAP.
        & Only needs \texttt{predict()}: permutation importance, KernelSHAP, LIME. \\
        \bottomrule
    \end{tabular}

    \vspace{1.1em}

    \raggedright
    \begin{block}{}
        \centering
        \small Model-agnostic buys generality at the price of speed and, often, fidelity.
        Model-specific is faster and usually more faithful, but locks you to a model family.
    \end{block}
\end{frame}

\begin{frame}{Try an Interpretable Model First}

    \begin{block}{Rudin (2019)}
        For high-stakes decisions, prefer a model that is interpretable by construction over a
        black box plus a post-hoc explanation --- because the explanation is an
        \textbf{approximation} of the model, and an approximation can be wrong in exactly the cases
        you care about.
    \end{block}

    \vspace{0.8em}

    \begin{itemize}
        \setlength\itemsep{0.45em}
        \item On many tabular problems, a well-regularised logistic regression, a depth-4 tree, or
              a generalised additive model is within a point or two of a gradient-boosted ensemble.
              \textbf{Measure the gap before assuming you need the black box.}
        \item If the gap is real and worth it, take the black box --- but then the explanation
              needs its own validation, exactly as the model does.
        \item Interpretability is not a property of the model alone. A linear model with 4{,}000
              engineered features is not interpretable; a depth-3 tree over four meaningful
              features is.
    \end{itemize}

    \vspace{0.3em}
    {\scriptsize C.~Rudin, ``Stop explaining black box machine learning models for high stakes
    decisions and use interpretable models instead,'' \emph{Nature Machine Intelligence} 1:206--215,
    2019, \href{https://arxiv.org/abs/1811.10154v3}{arXiv:1811.10154v3}.}
\end{frame}

